\section{Related Work}

\subsection{Low-Rank Matrix Factorization}

\textbf{Eckart--Young theorem} (Eckart \& Young, 1936) establishes that truncated 
SVD is optimal in Frobenius norm at fixed rank. Our method uses this as a 
building block but \textbf{selects rank adaptively} via energy thresholding.

\textbf{Randomized SVD} (Halko et al., 2011) accelerates the decomposition. We 
use exact SVD because our matrices are modest size ($768 \times 768$).

\textbf{Energy-based rank selection} is used in PCA (Jolliffe, 2002) and latent 
semantic analysis (Deerwester et al., 1990). We apply it to attention 
projections for the first time.

\subsection{Attention Compression}

\textbf{SOA Fused Attention} (companion work) uses a \textbf{fixed} rank $k$ for 
Q/K/V. We differ in two ways: (a) dynamic rank per projection, (b) all four 
projections.

\textbf{A3} and \textbf{AASVD} (companion works) use activation-weighted SVD requiring 
calibration data. We are \textbf{zero-shot}.

\textbf{Linformer} (Wang et al., 2020) projects K, V to a fixed rank at the 
sequence level. Our projection is at the feature level.

\subsection{Parameter-Efficient Fine-Tuning}

\textbf{LoRA} (Hu et al., 2021) adds trainable low-rank deltas to frozen weights. 
Our method adds LoRA on top of factorized SVD -- a two-level low-rank 
structure.

\textbf{AdaLoRA} (Zhang et al., 2023) also allocates rank adaptively, but for the 
adapter, not for the base factorization.

\subsection{T5 Relative Attention}

\textbf{T5} (Raffel et al., 2020) uses a learned relative position bias added to 
attention scores. This mechanism is orthogonal to our factorization but 
requires special handling in the weight transfer.

\subsection{Cross-Attention}

Encoder-decoder models (T5, BART) have cross-attention with separate K, V 
projections from the encoder. Our method supports cross-attention because 
all projections are factorized.

\subsection{Position of Our Work}

\begin{table}[H]
\centering
\small
\begin{tabular}{llllll}
\toprule
\textbf{Method} & \textbf{Rank} & \textbf{Calibration} & \textbf{Projections} & \textbf{T5} & \textbf{Cross-attn} \\
\midrule
SOA Fused & Fixed & No & Q/K/V & No & No \\
A3 & Fixed & Yes & Q/K/V/O & No & No \\
AASVD & Fixed & Yes & Q/K/V/O & No & No \\
\textbf{Adaptive SVD} & \textbf{Dynamic} & \textbf{No} & \textbf{Q/K/V/O} & \textbf{Yes} & \textbf{Yes} \\
\bottomrule
\end{tabular}
\end{table}

\textbf{Key distinction.} Adaptive SVD Attention is the only method that:
\begin{itemize}
\item Uses \textbf{dynamic rank per projection}
\item Is \textbf{zero-shot} (no calibration)
\item Supports \textbf{T5 relative attention}
\item Handles \textbf{encoder-decoder cross-attention}

\end{itemize}


\subsection{Comparison Table}

\begin{lstlisting}

---
\end{lstlisting}